% 宇宙学（综述）
% license CCBYSA3
% type Wiki

本文根据 CC-BY-SA 协议转载翻译自维基百科\href{https://en.wikipedia.org/wiki/Cosmology}{相关文章}。

宇宙学（Cosmology，源自古希腊语 κόσμος（cosmos，意为“宇宙”或“世界”）与 λογία（logia，意为“研究”））是一门研究宇宙——即宇宙整体本质的物理学与形而上学的分支。英语单词 “cosmology” 首次出现在 1656 年托马斯·布朗特（Thomas Blount）的《Glossographia》中，其含义为“关于世界的论述”。\(^\text{[2]}\)1731 年，德国哲学家克里斯蒂安·沃尔夫（Christian Wolff）在拉丁语中使用 “cosmologia” 一词，用以指代研究物质世界一般性质的形而上学分支。\(^\text{[3]}\)宗教或神话宇宙学（religious or mythological cosmology）是一套基于神话、宗教与秘传文献及传统的信仰体系，通常与创世神话（creation myths）及末世论（eschatology）相关。在天文学的科学范畴中，宇宙学主要关注对宇宙年代学（chronology of the universe）的研究。

物理宇宙学（physical cosmology）研究可观测宇宙的起源、其大尺度结构与动力学特征，以及宇宙的最终命运，同时探讨支配这些领域的科学定律。\(^\text{[4]}\)这一学科由科学家（如天文学家与物理学家）以及哲学家（如形而上学家、物理哲学家与时空哲学家）共同研究。由于其研究范围与哲学存在交叉，物理宇宙学中的理论既可能包含科学命题，也可能包含非科学命题，且可能依赖于无法通过实验检验的假设。物理宇宙学是天文学的一个子分支，专注于研究整个宇宙。现代物理宇宙学由“大爆炸理论”（Big Bang Theory）主导，该理论试图结合观测天文学与粒子物理学。\(^\text{[5][6]}\)更具体而言，标准的宇宙参数化模型包括暗物质（dark matter）与暗能量（dark energy），统称为“ΛCDM 模型”（Lambda-CDM model）。

理论天体物理学家戴维·N·斯珀格尔（David N. Spergel）将宇宙学称为一门“历史科学”（historical science），因为“当我们仰望宇宙时，实际上是在回望过去”，这源于光速的有限性。\(^\text{[7]}\)

